\documentclass[11pt,a4paper]{article}

% --- Encoding and fonts ---
\usepackage[utf8]{inputenc}
\usepackage[T1]{fontenc}

% --- Mathematics ---
\usepackage{amsmath,amssymb,amsthm}
\usepackage{mathtools}

% --- Page geometry ---
\usepackage[margin=1in,headheight=14pt]{geometry}

% --- Hyperlinks ---
\usepackage[colorlinks=true,linkcolor=red,citecolor=blue,urlcolor=blue]{hyperref}

% --- Lists ---
\usepackage{enumitem}

% --- Colors and boxes ---
\usepackage{xcolor}
\usepackage[theorems,skins,breakable]{tcolorbox}

% --- Header/Footer ---
\usepackage{fancyhdr}
\pagestyle{fancy}
\fancyhf{}
\fancyhead[L]{\textsc{Lesson 8: Convex Analysis}}
\fancyhead[R]{\thepage}
\renewcommand{\headrulewidth}{0.4pt}

% --- Tables ---
\usepackage{booktabs}

% --------------------------------------------------------------------------
% Theorem environments
% --------------------------------------------------------------------------
\newtheorem{theorem}{Theorem}[section]
\newtheorem{lemma}[theorem]{Lemma}
\newtheorem{proposition}[theorem]{Proposition}
\newtheorem{corollary}[theorem]{Corollary}

\theoremstyle{definition}
\newtheorem{definition}[theorem]{Definition}
\newtheorem{example}[theorem]{Example}
\newtheorem{exercise}[theorem]{Exercise}

\theoremstyle{remark}
\newtheorem{remark}[theorem]{Remark}

% --------------------------------------------------------------------------
% Colored boxes
% --------------------------------------------------------------------------
\newtcolorbox{keyresult}[1][]{%
  colback=green!5!white,
  colframe=green!50!black,
  fonttitle=\bfseries,
  title={Key Result},
  breakable,
  #1
}

\newtcolorbox{rlconnection}[1][]{%
  colback=blue!5!white,
  colframe=blue!50!black,
  fonttitle=\bfseries,
  title={RL/ML Connection},
  breakable,
  #1
}

\newtcolorbox{warning}[1][]{%
  colback=red!5!white,
  colframe=red!50!black,
  fonttitle=\bfseries,
  title={Warning},
  breakable,
  #1
}

% --------------------------------------------------------------------------
% Custom commands
% --------------------------------------------------------------------------
\newcommand{\R}{\mathbb{R}}
\newcommand{\N}{\mathbb{N}}
\newcommand{\Q}{\mathbb{Q}}
\newcommand{\Z}{\mathbb{Z}}
\newcommand{\C}{\mathbb{C}}
\newcommand{\Prob}{\mathbb{P}}
\newcommand{\E}{\mathbb{E}}
\newcommand{\indicator}{\mathbf{1}}

\DeclareMathOperator{\dom}{dom}
\DeclareMathOperator{\epi}{epi}
\DeclareMathOperator{\conv}{conv}
\DeclareMathOperator{\relint}{relint}
\DeclareMathOperator{\aff}{aff}
\DeclareMathOperator{\interior}{int}
\DeclareMathOperator{\cl}{cl}
\DeclareMathOperator{\tr}{tr}
\DeclareMathOperator{\diag}{diag}
\DeclareMathOperator{\rank}{rank}
\DeclareMathOperator{\sgn}{sgn}

% --------------------------------------------------------------------------
\begin{document}
\maketitle
\tableofcontents
\newpage

\setcounter{section}{0}

\section{Introduction: Why Convex Analysis for RL/ML}

Convex analysis provides the geometric and analytic foundation for optimization, which
in turn underlies virtually every learning algorithm. This document develops the theory
of convex sets and convex functions from first principles, with complete proofs,
following the treatment in Boyd \& Vandenberghe Chapters 2--3.

\textbf{What we cover.}
\begin{enumerate}[label=(\roman*)]
  \item \textbf{Convex sets}: definitions, key examples, operations preserving convexity,
        and the supporting/separating hyperplane theorems.
  \item \textbf{Convex functions}: first- and second-order characterizations with full
        proofs, Jensen's inequality, and the epigraph connection.
  \item \textbf{Operations preserving convexity} of functions: pointwise suprema,
        composition rules, and perspective transforms.
  \item \textbf{Conjugate functions}: the Fenchel conjugate, worked examples, and
        Fenchel's inequality.
  \item \textbf{Sublevel sets and quasiconvexity}: definitions and basic properties.
  \item \textbf{Log-concave and log-convex functions}: closure properties with
        applications to probability.
\end{enumerate}

\begin{rlconnection}
Convexity appears throughout machine learning and reinforcement learning:
\begin{itemize}
  \item \textbf{Loss functions}: Most supervised learning objectives (squared error,
        cross-entropy, hinge loss) are convex in the model parameters for linear models.
  \item \textbf{Policy optimization}: The KL-divergence constraint in TRPO/PPO defines a
        convex trust region. Understanding convex duality clarifies why the constrained
        policy update admits an efficient solution.
  \item \textbf{Value function geometry}: The Bellman optimality operator preserves
        concavity of the value function in certain state-action structures.
  \item \textbf{Fenchel duality}: Conjugate functions connect primal and dual
        optimization formulations, central to regularized MDPs and inverse RL.
  \item \textbf{Jensen's inequality}: Foundational for bounding expectations, variance
        reduction, and importance sampling analysis.
\end{itemize}
\end{rlconnection}

\textbf{Prerequisites.} Familiarity with real analysis (limits, continuity,
differentiability), linear algebra (eigenvalues, positive semidefiniteness), and basic
probability.

%% ========================================================================
%% SECTION 2: CONVEX SETS
%% ========================================================================

\section{Convex Sets}

\subsection{Basic Definitions}

\begin{definition}[Affine Set]
\label{def:affine-set}
A set $C \subseteq \R^n$ is \emph{affine} if for every $x_1, x_2 \in C$ and every
$\theta \in \R$, the point $\theta x_1 + (1-\theta) x_2 \in C$. Equivalently, $C$
contains the entire line through any two of its points.
\end{definition}

\begin{definition}[Convex Set]
\label{def:convex-set}
A set $C \subseteq \R^n$ is \emph{convex} if for every $x_1, x_2 \in C$ and every
$\theta \in [0,1]$, we have
\[
  \theta x_1 + (1-\theta) x_2 \in C.
\]
\end{definition}

\begin{definition}[Convex Combination and Convex Hull]
\label{def:convex-hull}
A \emph{convex combination} of points $x_1, \ldots, x_k$ is any point of the form
$\sum_{i=1}^k \theta_i x_i$ where $\theta_i \ge 0$ and $\sum_{i=1}^k \theta_i = 1$.
The \emph{convex hull} of a set $S$ is
\[
  \conv(S) = \left\{ \sum_{i=1}^k \theta_i x_i \;\middle|\;
    x_i \in S, \; \theta_i \ge 0, \;
    \sum_{i=1}^k \theta_i = 1, \; k \in \N \right\}.
\]
\end{definition}

\begin{proposition}[Convex Hull is the Smallest Convex Set]
\label{prop:convex-hull-smallest}
For any $S \subseteq \R^n$, $\conv(S)$ is convex, and $\conv(S) \subseteq C$ for every
convex set $C$ containing $S$.
\end{proposition}

\begin{proof}
Let $x = \sum_{i=1}^k \alpha_i x_i$ and $y = \sum_{j=1}^m \beta_j y_j$ be two elements
of $\conv(S)$, with $\alpha_i, \beta_j \ge 0$, $\sum \alpha_i = 1$, $\sum \beta_j = 1$.
For $\theta \in [0,1]$,
\[
  \theta x + (1-\theta) y = \sum_{i=1}^k \theta \alpha_i x_i
    + \sum_{j=1}^m (1-\theta) \beta_j y_j.
\]
The coefficients $\theta \alpha_i$ and $(1-\theta)\beta_j$ are nonnegative and sum to
$\theta + (1-\theta) = 1$, so this is a convex combination of points in $S$. Hence
$\theta x + (1-\theta)y \in \conv(S)$, proving convexity.

Now let $C$ be any convex set with $S \subseteq C$. By induction on $k$, every convex
combination of points in $S$ lies in $C$ (the base case $k=1$ is immediate; for the
inductive step, write $\sum_{i=1}^{k+1} \theta_i x_i = \theta_{k+1} x_{k+1} +
(1-\theta_{k+1}) \sum_{i=1}^k \frac{\theta_i}{1-\theta_{k+1}} x_i$, which is a convex
combination of $x_{k+1} \in C$ and a point in $C$ by the inductive hypothesis). Thus
$\conv(S) \subseteq C$.
\end{proof}

\begin{definition}[Cone]
\label{def:cone}
A set $C$ is a \emph{cone} if for every $x \in C$ and $\theta \ge 0$, we have
$\theta x \in C$. A set that is both convex and a cone is called a \emph{convex cone}.
Equivalently, $C$ is a convex cone if and only if for all $x_1, x_2 \in C$ and
$\theta_1, \theta_2 \ge 0$, we have $\theta_1 x_1 + \theta_2 x_2 \in C$.
\end{definition}

\subsection{Important Examples}

\begin{example}[Hyperplanes and Halfspaces]
\label{ex:hyperplane}
A \emph{hyperplane} is a set $\{x \in \R^n : a^T x = b\}$ for some $a \ne 0$. It is
affine (hence convex). A \emph{halfspace} is $\{x : a^T x \le b\}$, which is convex.
\end{example}

\begin{example}[Euclidean Balls and Ellipsoids]
\label{ex:ball-ellipsoid}
The Euclidean ball $B(x_c, r) = \{x : \|x - x_c\|_2 \le r\}$ is convex. More generally,
an \emph{ellipsoid} is $\mathcal{E} = \{x : (x - x_c)^T P^{-1} (x - x_c) \le 1\}$
where $P \succ 0$. Both are convex since the $\ell_2$-norm (and any norm) is a convex
function, and sublevel sets of convex functions are convex (Proposition~\ref{prop:sublevel}).
\end{example}

\begin{example}[Polyhedra]
\label{ex:polyhedra}
A \emph{polyhedron} is the intersection of finitely many halfspaces and hyperplanes:
$\mathcal{P} = \{x : Ax \le b, \; Cx = d\}$. Every polyhedron is convex.
\end{example}

\begin{example}[Positive Semidefinite Cone]
\label{ex:psd-cone}
The set $\mathbb{S}^n_+ = \{X \in \R^{n \times n} : X = X^T, \; X \succeq 0\}$ is a
convex cone. Indeed, if $A, B \succeq 0$ and $\theta_1, \theta_2 \ge 0$, then for any
$v \in \R^n$:
\[
  v^T(\theta_1 A + \theta_2 B)v = \theta_1 v^T A v + \theta_2 v^T B v \ge 0.
\]
\end{example}

\begin{rlconnection}
The PSD cone $\mathbb{S}^n_+$ appears in covariance matrix estimation (the covariance
matrix is always PSD) and in semidefinite relaxations used for robust control and
multi-agent coordination problems.
\end{rlconnection}

\subsection{Operations Preserving Convexity of Sets}

\begin{proposition}[Intersection]
\label{prop:intersection}
If $\{C_\alpha\}_{\alpha \in A}$ is any (possibly infinite) family of convex sets, then
$\bigcap_{\alpha \in A} C_\alpha$ is convex.
\end{proposition}

\begin{proof}
Let $x, y \in \bigcap_\alpha C_\alpha$ and $\theta \in [0,1]$. Then $x, y \in C_\alpha$
for every $\alpha$, so $\theta x + (1-\theta) y \in C_\alpha$ for every $\alpha$ (by
convexity of each $C_\alpha$). Hence $\theta x + (1-\theta) y \in \bigcap_\alpha
C_\alpha$.
\end{proof}

\begin{proposition}[Affine Image and Preimage]
\label{prop:affine-image}
Let $f(x) = Ax + b$ be an affine map. If $C$ is convex, then $f(C) = \{Ax + b : x \in
C\}$ is convex. If $D$ is convex, then $f^{-1}(D) = \{x : Ax + b \in D\}$ is convex.
\end{proposition}

\begin{proof}
For the image: let $u = Ax_1 + b, v = Ax_2 + b \in f(C)$ with $x_1, x_2 \in C$. Then
\[
  \theta u + (1-\theta)v = A(\theta x_1 + (1-\theta) x_2) + b = f(\theta x_1 +
  (1-\theta) x_2) \in f(C),
\]
since $\theta x_1 + (1-\theta) x_2 \in C$.

For the preimage: let $x_1, x_2 \in f^{-1}(D)$, so $Ax_1 + b, Ax_2 + b \in D$. Then
$A(\theta x_1 + (1-\theta) x_2) + b = \theta(Ax_1 + b) + (1-\theta)(Ax_2 + b) \in D$
by convexity of $D$, so $\theta x_1 + (1-\theta) x_2 \in f^{-1}(D)$.
\end{proof}

\begin{proposition}[Perspective Function]
\label{prop:perspective}
The \emph{perspective function} $P : \R^{n+1} \to \R^n$ defined by $P(x, t) = x/t$ with
$\dom P = \{(x,t) : t > 0\}$ preserves convexity: if $C \subseteq \dom P$ is convex,
then $P(C)$ is convex; if $D \subseteq \R^n$ is convex, then $P^{-1}(D) = \{(x,t) :
x/t \in D, \; t > 0\}$ is convex.
\end{proposition}

\begin{proof}
For the image: let $(x_1, t_1), (x_2, t_2) \in C$ with $t_1, t_2 > 0$ and
$\theta \in [0,1]$. The convex combination is $(\bar{x}, \bar{t}) = (\theta x_1 +
(1-\theta) x_2, \; \theta t_1 + (1-\theta) t_2)$, and $\bar{t} > 0$. Then
\[
  P(\bar{x}, \bar{t}) = \frac{\theta x_1 + (1-\theta) x_2}{\theta t_1 + (1-\theta) t_2}
  = \frac{\theta t_1}{\bar{t}} \cdot \frac{x_1}{t_1}
    + \frac{(1-\theta) t_2}{\bar{t}} \cdot \frac{x_2}{t_2}.
\]
The coefficients $\theta t_1 / \bar{t}$ and $(1-\theta) t_2 / \bar{t}$ are nonnegative and
sum to 1, so this is a convex combination of $P(x_1, t_1)$ and $P(x_2, t_2)$, both in
$P(C)$. Hence $P(\bar{x}, \bar{t}) \in P(C)$.

For the preimage: let $(x_1, t_1), (x_2, t_2) \in P^{-1}(D)$. Then $x_1/t_1, x_2/t_2
\in D$. The same algebra shows $P(\theta(x_1,t_1) + (1-\theta)(x_2,t_2))$ is a convex
combination of $x_1/t_1$ and $x_2/t_2$, hence in $D$.
\end{proof}

\subsection{Separating and Supporting Hyperplane Theorems}

\begin{theorem}[Separating Hyperplane Theorem]
\label{thm:separating}
Let $C$ and $D$ be nonempty, disjoint, convex subsets of $\R^n$. Then there exist
$a \in \R^n$, $a \ne 0$, and $b \in \R$ such that
\[
  a^T x \le b \quad \text{for all } x \in C, \qquad
  a^T x \ge b \quad \text{for all } x \in D.
\]
\end{theorem}

\begin{proof}
We prove the case where $C$ and $D$ are closed and at least one is compact (the general
case requires the Hahn--Banach theorem, which belongs to functional analysis).

Since $C$ and $D$ are closed with $C$ compact and $C \cap D = \emptyset$, the function
$\phi(x, y) = \|x - y\|_2$ attains its minimum on $C \times D$ at some pair
$(x^*, y^*)$. (Compactness of $C$ and closedness of $D$ together with the fact that
$\|x - y\| \to \infty$ as $\|y\| \to \infty$ guarantee the minimum is attained.) Let
$d = \|x^* - y^*\|_2 > 0$.

Set $a = y^* - x^*$ and $b = \tfrac{1}{2}(a^T y^* + a^T x^*) =
\tfrac{1}{2}(\|y^*\|^2 - \|x^*\|^2)$ (after expanding $a = y^* - x^*$). We claim
$a^T x \le b$ for all $x \in C$ and $a^T y \ge b$ for all $y \in D$.

For $x \in C$: consider $g(t) = \|x^* + t(x - x^*) - y^*\|_2^2$ for $t \in [0,1]$.
Since $x^* + t(x - x^*) \in C$ and $(x^*, y^*)$ minimizes the distance, $g(t) \ge
g(0)$ for all $t \ge 0$, so $g'(0) \ge 0$:
\[
  g'(0) = 2(x - x^*)^T(x^* - y^*) \ge 0,
\]
i.e., $(x - x^*)^T a \le 0$, giving $a^T x \le a^T x^*$.

Similarly, for $y \in D$: $h(t) = \|x^* - y^* - t(y - y^*)\|_2^2 \ge h(0)$
gives $h'(0) \ge 0$, i.e., $(y - y^*)^T(y^* - x^*) \ge 0$, so $a^T y \ge a^T y^*$.

Since $a^T x \le a^T x^*$ for all $x \in C$ and $a^T y \ge a^T y^*$ for all $y \in D$,
and $a^T x^* < a^T y^*$ (because $a^T(y^* - x^*) = \|a\|^2 > 0$), any
$b \in [a^T x^*, a^T y^*]$ works. We choose $b = (a^T x^* + a^T y^*)/2$.
\end{proof}

\begin{theorem}[Supporting Hyperplane Theorem]
\label{thm:supporting}
Let $C \subseteq \R^n$ be a nonempty convex set and let $x_0 \in \partial C$ (the
boundary of $C$). Then there exists $a \ne 0$ such that $a^T x \le a^T x_0$ for all
$x \in C$. The hyperplane $\{x : a^T x = a^T x_0\}$ is called a \emph{supporting
hyperplane} to $C$ at $x_0$.
\end{theorem}

\begin{proof}
Since $x_0 \in \partial C$, there is a sequence $\{y_k\}$ with $y_k \notin C$ and
$y_k \to x_0$. For each $k$, apply the separating hyperplane theorem to the convex sets
$C$ and $\{y_k\}$: there exists $a_k$ with $\|a_k\| = 1$ and $a_k^T x \le a_k^T y_k$
for all $x \in C$. Since $\{a_k\}$ lies on the unit sphere (compact), it has a
convergent subsequence $a_{k_j} \to a$ with $\|a\| = 1$. Passing to the limit:
$a^T x \le a^T x_0$ for all $x \in C$.
\end{proof}

\begin{rlconnection}
The separating hyperplane theorem is the geometric engine behind Lagrangian duality.
Every constraint in an optimization problem can be interpreted as requiring the
decision variable to lie in a convex set, and Lagrange multipliers define the
separating hyperplane. In constrained policy optimization (CPO), the safety constraint
defines a convex feasible region in surrogate-objective space, and the dual variable
(Lagrange multiplier) encodes the ``price'' of constraint violation.
\end{rlconnection}

%% ========================================================================
%% SECTION 3: CONVEX FUNCTIONS
%% ========================================================================


\end{document}